\section{Future: Evaluation, Multimodal RAG y RAG 2.0}

La última parte cierra el curso con una agenda de investigación sobre sistemas: el joint pretraining desde cero sigue sin ser suficiente; la scaling law, el chunker, el encoder, la database, los synthetic data y la evaluation pueden aprenderse o rediseñarse; RAG también puede ir más allá del texto, hacia imágenes y la multimodalidad; el objetivo final es optimizar todo el cost--quality system.

\subsection{Open Questions: qué capas todavía no se han optimizado}

La slide de Open Questions plantea preguntas sobre el joint pretraining, las scaling laws, el reranker/chunker, las representaciones más allá del bi-encoder, la database convergence, los synthetic data y la evaluation. Muestra que la idea de que «RAG ya está resuelto» dista mucho de ser un hecho: los sistemas actuales todavía dependen en gran medida de la segmentación manual y de métricas indirectas débiles.

\lecturefigure{slide-47-open-questions.jpg}{Problemas abiertos de RAG en joint training, chunking, database y evaluation}{Video de la clase de Stanford 01:07:36}

\begin{knowledgebox}{Matriz por capas de la RAG Evaluation}
\begin{tabularx}{\textwidth}{L{0.18\textwidth}>{\raggedright\arraybackslash}X >{\raggedright\arraybackslash}X}
\toprule
Capa & Ejemplos de métricas & Punto ciego típico \\
\midrule
Retrieval & Recall@$k$, MRR, nDCG & la relevance label no equivale a la generator utility \\
Packing & evidence coverage, position & ignora el duplicate y el token budget \\
Generation & correctness, citation, faithfulness & una respuesta correcta puede provenir de la memoria paramétrica \\
System & latency, cost, freshness, access control & el promedio oculta la latencia de cola y los errores de alto riesgo \\
\bottomrule
\end{tabularx}
\end{knowledgebox}

\teachervoice{El ponente critica que la RAG evaluation actual es débil: la downstream performance no permite ver dónde falla la retrieval, y calcular la harmonic mean de la retrieval accuracy y la LM accuracy tampoco cuenta con el respaldo de conjuntos de datos confiables.}

\subsection{¿Se integrará la Vector Database en la Database convencional?}

Entre las preguntas abiertas que planteó de forma oral, el ponente duda de que la vector database especializada vaya a existir de manera independiente a largo plazo: BM25 es más eficiente y el reranker puede complementar la parte semántica; las bases de datos tradicionales están incorporando capacidades de recuperación vectorial, y los proveedores de bases de datos vectoriales también están añadiendo sparse retrieval y metadata filtering. Se trata de un juicio sobre la industria de 2023; debe discutirse como una tendencia de arquitectura y no tomarse como un hecho del mercado actual.

\begin{warningbox}{No conviertas las predicciones de la clase en conclusiones sobre productos}
La implementación de las bases de datos, el hardware, los algoritmos de indexación y las capacidades de los proveedores cambian con rapidez. Las notas conservan solo la pregunta abstracta: si los sistemas futuros necesitarán gestionar de forma unificada sparse, dense, metadata, transactions y access control, sin afirmar que algún tipo de producto haya desaparecido o haya ganado.
\end{warningbox}

\subsection{Multimodal RAG}

La slide de Multimodal muestra la cross-modal retrieval augmentation y un pipeline de visual question answering. Una query de texto puede recuperar imágenes, una imagen puede recuperar texto, o bien un vision pipeline genera primero una descripción estructurada y después entrega el resultado a un frozen LM.

\lecturefigure{slide-48-multimodal-rag.jpg}{Cross-modal retrieval y multimodal retrieval-augmented LM}{Video de la clase de Stanford 01:08:45}

\begin{knowledgebox}{Lectura de la figura: el Multimodal RAG necesita alinear tres espacios}
El sistema debe alinear, como mínimo, la query representation, la retrieved modality representation y la generator input representation. Que dos imágenes sean similares no significa que sean pertinentes para la respuesta, y la caption también puede perder detalles visuales; por eso se necesitan un cross-modal retriever, un modality adapter y una task-level evaluation.
\end{knowledgebox}

\begin{warningbox}{Los resultados de la recuperación visual también pueden convertirse en evidencia errónea}
Las imágenes similares pueden provenir de momentos, lugares u objetos distintos, y el OCR/caption puede equivocarse. El Multimodal grounding requiere provenance y verificación de fechas y entidades; no se debe elevar el nivel de la evidencia solo porque «hay una imagen».
\end{warningbox}

\teachervoice{Kiela considera la multimodality una extensión natural: ¿por qué RAG debería quedarse solo en el text? Cita LENS y el cross-modal work, y enfatiza que incluso un frozen LM puede adquirir capacidades visuales mediante un vision pipeline externo.}

\subsection{RAG 2.0: Systems over Models}

La slide final cierra con cuatro puntos: systems over models, optimize it all, trade cost and quality y zero-shot domain generalization. RAG 2.0 no es un diagrama fijo, sino una optimization stance: retriever, chunker, reranker, generator, index, data y serving no deberían separarse de manera permanente.

\lecturefigure{slide-49-rag-2-0.jpg}{RAG 2.0: de los frozen components a systems-over-models}{Video de la clase de Stanford 01:10:12}

\begin{importantbox}{El objetivo de sistema de Cost--Quality}
El objetivo de despliegue puede abstraerse como
\[
\max_{\Theta,I,\pi}\;
Q(\Theta,I,\pi)
-\lambda_L L(\Theta,I,\pi)
-\lambda_C C(\Theta,I,\pi)
-\lambda_R R(\Theta,I,\pi),
\]
donde $Q$ es la calidad de la tarea, $L$ es la latency, $C$ es el compute/storage cost y $R$ es el riesgo o las salidas no confiables; $\pi$ incluye la retrieval policy. Optimizar solo la LM accuracy no basta para determinar la solución óptima del sistema.
\end{importantbox}

\teachervoice{El ponente resume RAG 2.0 como ``systems over models'' e incluso sugiere retropropagar hasta el chunker. Lo esencial no es que todos los módulos deban actualizarse al mismo tiempo, sino que cualquier frontera definida a mano debe someterse a la prueba de «si puede aprenderse y si vale la pena optimizarla».}

\subsection{Closing Q\&A: Hybrid, Fine-tuning y Retrieval Hardware}

La Q\&A final ofrece tres juicios prácticos. Primero, BM25 puede servir para cast a wide net y después un dense reranker acota los resultados. Segundo, la domain adaptation terminará combinando retrieval, fine-tuning y adapters, en lugar de elegir siempre una de dos opciones. Tercero, el ponente menciona la posibilidad industrial de un dedicated retrieval hardware, pero no ofrece modelos ni benchmark verificables públicamente, por lo que solo se registra como speaker speculation.

\begin{knowledgebox}{Retrieval y Fine-tuning resuelven estados distintos}
El fine-tuning modifica los parameters y es adecuado para la conducta, el formato, las representaciones del dominio y las habilidades estables; la retrieval modifica la per-request evidence y es adecuada para actualizaciones frecuentes, datos privados y attribution. RAG 2.0 puede aplicar fine-tune a todo el retrieval-augmented system en tareas de dominio y, al mismo tiempo, conservar un index actualizable.
\end{knowledgebox}

\teachervoice{El ponente afirma explícitamente que, al final, probablemente todos combinarán retrieval, end-to-end fine-tuning y adapters. La verdadera pregunta pasa a ser cómo combinarlos de forma eficiente, y no discutir cuál es la única opción posible.}

\subsection{Hallucination, Generic Error y Ground Truth}

El pasaje spoken-only más importante del curso aparece en los últimos cuatro minutos. Al ponente le inconforma que el campo llame hallucination a todos los errores: los errores de cálculo comunes, la mala interpretación de la pregunta y la evidence contradiction deben distinguirse; la hallucination tiene que definirse, como mínimo, respecto de alguna ground truth o de la retrieved evidence.

\begin{importantbox}{Distinción entre tres tipos de fallas}
\begin{tabularx}{\textwidth}{L{0.20\textwidth}>{\raggedright\arraybackslash}X >{\raggedright\arraybackslash}X}
\toprule
Tipo & Criterio & Ejemplo \\
\midrule
Generic error & La salida no cumple con la respuesta de la tarea ni con las reglas de cálculo & Errores aritméticos, restricciones omitidas, errores de formato \\
Unsupported claim & No hay evidence suficiente que la respalde & El index no tiene fuentes pertinentes y aun así se da una afirmación categórica \\
Evidence contradiction & La salida entra en conflicto con la ground truth/source seleccionada & El documento dice A y la respuesta afirma no A \\
\bottomrule
\end{tabularx}
\end{importantbox}

\begin{warningbox}{La elección de la ground truth es en sí misma una decisión de sistema}
El index puede admitir solo peer-reviewed papers o incluir también arXiv, documentos empresariales o la web abierta; distintas source policy producen distintos «hechos confiables». RAG hace explícita esta elección, pero una source errónea o maliciosa todavía puede volver errónea una grounded answer.
\end{warningbox}

\teachervoice{El ponente señala que un frozen GPT-4, incluso conectado a RePlug, todavía puede ignorar la evidence; si todo el conocimiento se trasladara a un index externo y el modelo se encargara solo del lenguaje y del reasoning, en teoría sería más fácil forzar el grounding, pero se trata de un límite que marca una dirección y no de una garantía ya alcanzada.}

\teachervoice{Añade que la creative writing a veces necesita «hallucinate», mientras que una factual task requiere un grounding más fuerte; un sistema ideal debería tener una grounding knob. La temperature solo cambia la sampling distribution: incluso con temperatura baja el modelo puede producir errores de forma estable, por lo que no puede funcionar como grounding control.}

\subsection{Resumen de la sección}

El Future RAG necesita integrar evaluation, database, multimodal retrieval, chunking, instruction y serving. La conclusión final de RAG 2.0 es optimizar todo el system y especificar con claridad la source of truth, la grounding strength, el cost y el risk; la recuperación y el fine-tuning son herramientas complementarias, y la temperature no es un interruptor de factualidad.
